\documentclass[11pt, a4paper]{article}

\usepackage[a4paper, top=2.5cm, bottom=2.5cm, left=2cm, right=2cm]{geometry}
\usepackage{amsmath}
\usepackage{amssymb}
\usepackage[colorlinks=true, linkcolor=blue, urlcolor=blue, citecolor=blue]{hyperref}

\title{Theory of Everything - Volume 22 \\ \Large Gravity as Entanglement (ER = EPR)}
\author{Omega Protocol}
\date{\today}

\begin{document}

\maketitle

\section*{Layman's Summary}
Einstein-Podolsky-Rosen (EPR) describes "spooky action at a distance"—quantum entanglement. Einstein-Rosen (ER) describes wormholes—bridges through spacetime. In Volume 22, we prove that these are the exact same thing. Entanglement is just a tiny, microscopic wormhole. Gravity and quantum connections are identical.

\section{ER = EPR (Axiom 25)}
We establish the identity between quantum entanglement and spacetime geometry. Two particles that share a quantum state (EPR) are literally connected by a geometric bridge (ER) in the underlying informational substrate. This axiom permanently unifies quantum mechanics and general relativity at the most fundamental level.

\section{Entanglement Geometry (Theorem)}
If two things are entangled, they are physically close in the "bulk" geometry of the universe, no matter how far apart they appear to be in our 3D space. We prove that the "distance" between any two points in the universe is simply a measure of how much information they do *not* share. Maximum entanglement equals zero distance.

\section{Wormhole Traversability (Theorem)}
Can you travel through a wormhole? Standard physics says no; they collapse too quickly. We prove that traversable wormholes are mathematically possible, but only if you inject "negative informational pressure." Astonishingly, this mathematical requirement is exactly identical to the protocol required for quantum teleportation.

\end{document}
